% paper7_billion_edge.tex — arxiv submission (cs.DB)
% Scaling Knowledge Graph Federation to One Billion Edges on Commodity Hardware
\documentclass[11pt]{article}

% --- Packages ---
\usepackage[utf8]{inputenc}
\usepackage[T1]{fontenc}
\usepackage{lmodern}
\usepackage[margin=1in]{geometry}
\usepackage{amsmath,amssymb}
\usepackage{graphicx}
\usepackage{booktabs}
\usepackage{hyperref}
\usepackage{xcolor}
\usepackage{listings}
\usepackage{natbib}
\usepackage{float}
\usepackage{enumitem}
\usepackage{url}
\usepackage{microtype}
\usepackage{multirow}

% --- Hyperref setup ---
\hypersetup{
  colorlinks=true,
  linkcolor=blue!70!black,
  citecolor=blue!70!black,
  urlcolor=blue!70!black
}

% --- Listings: Cypher ---
\lstdefinelanguage{Cypher}{
  keywords={MATCH, RETURN, WHERE, CREATE, WITH, ORDER, BY, LIMIT, AS, DISTINCT, AND, OR, IN, OPTIONAL},
  sensitive=false,
  morestring=[b]',
  morestring=[b]",
  morecomment=[l]{//},
}
\lstset{
  language=Cypher,
  basicstyle=\ttfamily\small,
  keywordstyle=\bfseries\color{blue!70!black},
  stringstyle=\color{red!70!black},
  commentstyle=\itshape\color{gray},
  breaklines=true,
  frame=single,
  framesep=3pt,
  xleftmargin=6pt,
  xrightmargin=6pt,
  aboveskip=8pt,
  belowskip=8pt
}

% --- Title ---
\title{Scaling Knowledge Graph Federation to One Billion Edges\\on Commodity Hardware}

\author{
  Madhulatha Mandarapu\thanks{madhulatha@samyama.ai, ORCID: \url{https://orcid.org/0009-0005-2837-6725}} \and
  Sandeep Kunkunuru\thanks{sandeep@samyama.ai, ORCID: \url{https://orcid.org/0000-0002-8886-1846}}
}
\date{%
  VaidhyaMegha Private Limited, India\\[2pt]
  \url{https://samyama.ai/}\\[8pt]
  April 2026
}

\begin{document}
\maketitle

% ============================================================
% ABSTRACT
% ============================================================
\begin{abstract}
Cross-domain knowledge graph federation---querying from molecular biology to population health in a single language---requires loading billions of edges from heterogeneous open data sources into a unified graph store. Existing solutions either fragment data across multiple systems (requiring ETL pipelines), or demand expensive infrastructure (multi-node clusters, proprietary engines). We present a single-node approach using the Samyama graph database (Rust, OpenCypher, RESP protocol) that loads \textbf{74.3 million nodes and 1.07 billion edges} from four biomedical knowledge graphs on a single AWS r6a.8xlarge instance for \textbf{\$2.50 in spot compute}. A two-phase bulk loading strategy---node stubs (777 bytes/node) and edge stubs (52 bytes/edge)---achieves 13.6$\times$ memory reduction over standard ingestion, enabling billion-edge graphs on 256\,GB RAM.

We evaluate 140 benchmark queries across two federated trifectas: a \textbf{biomedical trifecta} (PubMed, Clinical Trials, Pathways, Drug Interactions; 96/100 queries pass) and a \textbf{public health trifecta} (Disease Surveillance, Health Determinants, Health Systems; 40/40 queries pass). The 6-KG federation bridges molecular biology to population health via shared properties (\texttt{Country.iso\_code}, \texttt{Drug.drugbank\_id}, \texttt{Gene.gene\_name}), answering questions like ``Which countries with high poverty also have low IHR emergency capacity and active clinical trials for tuberculosis?'' All data, loaders, snapshots, and benchmark queries are open-source.
\end{abstract}

\noindent\textbf{Keywords:} Knowledge Graphs, Cross-Domain Federation, Billion-Edge Scale, Biomedical Data Integration, Public Health, OpenCypher, Rust, Commodity Hardware.

% ============================================================
% 1. INTRODUCTION
% ============================================================
\section{Introduction}
\label{sec:intro}

Biomedical and public health research increasingly requires integrating data across domain boundaries: understanding how a drug's gene targets relate to disease outbreaks in specific populations, or correlating a country's health expenditure with its vaccine coverage and clinical trial activity. This data exists in open repositories---PubMed (35M+ articles), ClinicalTrials.gov (480K+ trials), WHO Global Health Observatory (190+ countries)---but is fragmented across incompatible formats, schemas, and access patterns.

Knowledge graphs provide a natural abstraction for integrating heterogeneous data~\citep{hogan2021knowledge}, and graph databases enable efficient traversal across these connections. However, loading billion-edge real-world datasets into a single queryable graph presents engineering challenges: memory management at scale, efficient bulk loading, and cross-domain federation without ETL between knowledge graphs.

We make four contributions:

\begin{enumerate}[leftmargin=*]
  \item \textbf{Billion-edge scale on commodity hardware}: 74.3M nodes and 1.07B edges loaded from 4 biomedical sources in 3 hours 46 minutes on a single AWS r6a.8xlarge instance (\$2.50 spot cost). Two-phase bulk loading with node stubs (777\,B/node) and edge stubs (52\,B/edge) achieves 13.6$\times$ memory reduction.

  \item \textbf{6-KG cross-domain federation}: Three biomedical KGs (PubMed, Clinical Trials, Drug Interactions) federate with three public health KGs (Disease Surveillance, Health Determinants, Health Systems) via shared bridge properties---no ETL, no data movement.

  \item \textbf{140-query benchmark}: 100 biomedical queries (96 pass) and 40 public health queries (40 pass) spanning point lookups, traversals, aggregations, and cross-KG federation. All queries, results, and timing published.

  \item \textbf{Full reproducibility}: All 8 KGs built from open data, all loaders open-source (Rust), all snapshots on S3 and GitHub, all queries in CSV format.
\end{enumerate}

% ============================================================
% 2. ARCHITECTURE
% ============================================================
\section{Architecture}
\label{sec:arch}

Samyama~\citep{mandarapu2026samyama} is a graph database written in Rust with OpenCypher query support, RESP wire protocol (Redis-compatible), and RocksDB persistence. For billion-edge loading, we extend the core architecture with three mechanisms:

\subsection{Two-Phase Bulk Loading}
\label{sec:bulk-loading}

Standard node creation in Samyama allocates a \texttt{HashMap} for properties, label index entries, and adjacency list initialization---3,785 bytes per node. At 66M nodes (PubMed), this requires 250\,GB for node allocation alone.

Two-phase bulk loading separates topology from properties:

\begin{enumerate}[leftmargin=*]
  \item \textbf{Phase 1: Stubs.} \texttt{create\_node\_stub(label)} allocates only a node ID, label, and empty adjacency entry (777\,bytes). \texttt{create\_edge\_stub(src, tgt, type)} adds a single adjacency entry (52\,bytes) without creating a full Edge object (709\,bytes in standard path). Memory reduction: \textbf{4.9$\times$} per node, \textbf{13.6$\times$} per edge.
  \item \textbf{Phase 2: Properties.} \texttt{set\_column\_property(node\_id, key, value)} writes directly to a sparse \texttt{ColumnStore} (HashMap-based per-label columns) without touching the per-node HashMap.
\end{enumerate}

\subsection{Hybrid CSR Adjacency}
\label{sec:csr}

After each loading phase, \texttt{compact\_adjacency()} freezes the current write buffer into a Compressed Sparse Row (CSR) segment. Multiple segments coexist: traversal iterates all segments, while new writes go to the unfrozen buffer. Mid-phase compaction (every 50M edges) prevents unbounded write-buffer growth. The PubMed graph uses 19 CSR segments for 1.04B edges.

\subsection{Memory Budget}
\label{sec:memory}

\begin{table}[H]
  \centering
  \caption{Memory breakdown for the 74.3M-node biomedical trifecta.}
  \label{tab:memory}
  \begin{tabular}{@{}lrrl@{}}
    \toprule
    \textbf{Component} & \textbf{Per-Unit} & \textbf{Total} & \textbf{Notes} \\
    \midrule
    Node stubs (66.2M)  & 777 B   & 51.4 GB  & PubMed articles \\
    Edge stubs (1.04B)  & 52 B    & 54.1 GB  & Adjacency entries \\
    ColumnStore         & varies  & ${\sim}$40 GB & String interning helps \\
    ID maps (dedup)     & varies  & ${\sim}$8 GB  & Dropped after each phase \\
    CSR frozen segments & compact & ${\sim}$70 GB & 19 segments \\
    \midrule
    \textbf{Peak RSS}   &         & \textbf{${\sim}$210 GB} & r6a.8xlarge (256 GB) \\
    \bottomrule
  \end{tabular}
\end{table}

% ============================================================
% 3. KNOWLEDGE GRAPH CATALOG
% ============================================================
\section{Knowledge Graph Catalog}
\label{sec:catalog}

We construct 8 knowledge graphs from 12 open data sources (Table~\ref{tab:catalog}). All loaders are Rust programs reading pre-downloaded JSON, CSV, and pipe-delimited files.

\begin{table}[H]
  \centering
  \caption{Knowledge graph catalog. All data is open-access.}
  \label{tab:catalog}
  \begin{tabular}{@{}p{3.2cm}rrp{4.5cm}@{}}
    \toprule
    \textbf{Knowledge Graph} & \textbf{Nodes} & \textbf{Edges} & \textbf{Data Sources} \\
    \midrule
    \multicolumn{4}{@{}l}{\textit{Biomedical Trifecta}} \\[2pt]
    PubMed/MEDLINE      & 66.2M  & 1.04B & NLM baseline files (XML) \\
    Clinical Trials      & 7.8M   & 27M   & AACT/ClinicalTrials.gov \\
    Pathways             & 119K   & 835K  & Reactome, STRING, Gene Ontology \\
    Drug Interactions    & 245K   & 388K  & DrugBank, SIDER, ChEMBL, DGIdb, OpenFDA \\
    \midrule
    \multicolumn{4}{@{}l}{\textit{Public Health Trifecta}} \\[2pt]
    Disease Surveillance & 217K   & 241K  & WHO GHO (diseases, vaccines, mortality) \\
    Health Determinants  & 286K   & 286K  & World Bank WDI (80 indicators), WHO Air Quality, WHO WASH \\
    Health Systems       & 20K    & 19K   & WHO SPAR v2 (IHR capacities), WHO NHWA (workforce) \\
    \midrule
    Cricket (baseline)   & 37K    & 1.4M  & Cricsheet.org \\
    \midrule
    \textbf{Total}       & \textbf{74.9M} & \textbf{1.07B} & \\
    \bottomrule
  \end{tabular}
\end{table}

\paragraph{Biomedical trifecta.} PubMed provides the citation graph (66.2M articles, 398M CITES edges, 342M AUTHORED\_BY, 342M ANNOTATED\_WITH). Clinical Trials adds 7.8M nodes (trials, conditions, interventions, sponsors). Pathways provides molecular biology (proteins, reactions, GO terms). Drug Interactions bridges pharmacology (drug-gene interactions, side effects, adverse events from 5 sources).

\paragraph{Public health trifecta.} Disease Surveillance covers WHO disease indicators, vaccine coverage, and mortality across 234 countries. Health Determinants provides 80 curated World Bank WDI indicators (socioeconomic, environmental, nutrition, demographic, water) for 211 countries from 1990--2024. Health Systems covers IHR emergency preparedness (15 SPAR v2 capacities) and health workforce density (physicians, nurses, dentists, pharmacists) from WHO.

% ============================================================
% 4. FEDERATION ARCHITECTURE
% ============================================================
\section{Federation Architecture}
\label{sec:federation}

Cross-KG queries require no ETL or data movement. Federation is achieved through \textbf{shared bridge properties}---common identifiers present in multiple KGs that enable Cypher MATCH clauses to traverse across domain boundaries.

\begin{table}[H]
  \centering
  \caption{Bridge properties for cross-KG federation.}
  \label{tab:bridges}
  \begin{tabular}{@{}llll@{}}
    \toprule
    \textbf{Bridge Property} & \textbf{From KG} & \textbf{To KG} & \textbf{Join Cardinality} \\
    \midrule
    \texttt{Country.iso\_code} & All public health & Clinical Trials & 211 countries \\
    \texttt{Drug.drugbank\_id} & Drug Interactions & Clinical Trials & ${\sim}$12K drugs \\
    \texttt{Gene.gene\_name}   & Drug Interactions & Pathways         & ${\sim}$20K genes \\
    \texttt{Protein.uniprot\_id} & Pathways        & Drug Interactions & ${\sim}$20K proteins \\
    \texttt{Article.nct\_id}   & PubMed           & Clinical Trials   & 1.02M articles \\
    \bottomrule
  \end{tabular}
\end{table}

\paragraph{Cross-trifecta example.} The following query spans all 6 KGs---from population vulnerability (Health Determinants) to emergency capacity (Health Systems) to disease burden (Surveillance) in a single Cypher session:

\begin{lstlisting}
// Which high-poverty countries have low IHR capacity?
MATCH (c1:Country)-[:HAS_INDICATOR]->(pov:SocioeconomicIndicator)
WHERE pov.indicator_code = 'SI.POV.DDAY' AND pov.value > 20
WITH c1.iso_code AS iso, pov.value AS poverty
ORDER BY poverty DESC LIMIT 5
MATCH (e:EmergencyResponse)-[:CAPACITY_FOR]->(c2:Country)
WHERE c2.iso_code = iso AND e.year = 2023
RETURN c2.name, poverty, e.capacity_code, e.score
\end{lstlisting}

This query pattern---property-value matching across separately loaded KGs---avoids the need for explicit foreign key relationships or pre-computed join tables. The cost-based query planner uses triple-level statistics to select efficient join orders.

% ============================================================
% 5. BENCHMARK DESIGN
% ============================================================
\section{Benchmark Design}
\label{sec:benchmark}

We design 140 queries across 7 categories (Table~\ref{tab:query-categories}).

\begin{table}[H]
  \centering
  \caption{Query categories and counts.}
  \label{tab:query-categories}
  \begin{tabular}{@{}llrrr@{}}
    \toprule
    \textbf{Category} & \textbf{Description} & \textbf{Biomed} & \textbf{PH} & \textbf{Total} \\
    \midrule
    Point lookup     & Single-node retrieval by ID       & 10 & 4  & 14 \\
    1-hop traversal  & Single relationship traversal     & 20 & 8  & 28 \\
    Aggregation      & GROUP BY / ORDER BY / LIMIT       & 25 & 10 & 35 \\
    Multi-hop        & 2+ relationship traversals        & 15 & 4  & 19 \\
    Cross-indicator  & Join across indicator categories   & 10 & 8  & 18 \\
    Cross-KG         & Federation across KG boundaries    & 15 & 6  & 21 \\
    Scan/filter      & Label scan with property filter    & 5  & 0  & 5 \\
    \midrule
    \textbf{Total}   &                                   & \textbf{100} & \textbf{40} & \textbf{140} \\
    \bottomrule
  \end{tabular}
\end{table}

All queries are published as CSV files with columns: ID, name, category, hop count, and Cypher query text. Results include timing (milliseconds), row count, and pass/empty/error status.

% ============================================================
% 6. EVALUATION
% ============================================================
\section{Evaluation}
\label{sec:eval}

\subsection{Biomedical Trifecta (74.3M nodes, 1.07B edges)}

\paragraph{Loading.} The PubMed graph loads in 6 phases on AWS r6a.8xlarge (32 vCPU, 256\,GB RAM, gp3 SSD):

\begin{table}[H]
  \centering
  \caption{PubMed loading phases (r6a.8xlarge, spot).}
  \label{tab:loading}
  \begin{tabular}{@{}lrrr@{}}
    \toprule
    \textbf{Phase} & \textbf{Nodes} & \textbf{Edges} & \textbf{Time} \\
    \midrule
    1. Articles      & 66.2M &  ---   & 42 min \\
    2. Authors       & 10M   & 342M   & 55 min \\
    3. MeSH terms    & 30K   & 342M   & 48 min \\
    4. Chemicals     & 500K  & 51M    & 12 min \\
    5. Citations     & ---   & 398M   & 52 min \\
    6. Grants        & 1M    & 28M    & 7 min \\
    \midrule
    \textbf{Total}   & \textbf{66.2M} & \textbf{1.04B} & \textbf{3h 46m} \\
    \bottomrule
  \end{tabular}
\end{table}

Clinical Trials, Pathways, and Drug Interactions add 8.1M nodes and 28M edges in under 5 minutes (standard \texttt{create\_node()} path---these are small enough for in-memory loading without stubs).

\paragraph{Cost.} AWS r6a.8xlarge spot price in ap-south-1 (Mumbai): \$0.67/hr $\times$ 3.77 hours = \textbf{\$2.52}. For comparison: Neo4j Aura Professional at this scale exceeds \$2,000/month; AWS Neptune would cost \$1,500+/month.

\paragraph{Querying.}

\begin{table}[H]
  \centering
  \caption{Biomedical benchmark results (100 queries).}
  \label{tab:biomed-results}
  \begin{tabular}{@{}lrrrr@{}}
    \toprule
    \textbf{Category} & \textbf{Queries} & \textbf{Pass} & \textbf{Empty} & \textbf{Timeout} \\
    \midrule
    PubMed            & 35 & 33 & 1 & 1 \\
    Clinical Trials   & 20 & 19 & 1 & 0 \\
    Pathways          & 15 & 15 & 0 & 0 \\
    Drug Interactions & 15 & 14 & 1 & 0 \\
    Cross-KG          & 15 & 15 & 0 & 0 \\
    \midrule
    \textbf{Total}    & \textbf{100} & \textbf{96} & \textbf{3} & \textbf{1} \\
    \bottomrule
  \end{tabular}
\end{table}

The 3 empty results are data gaps (specific PMIDs without grant data, missing edge types in snapshots). The 1 timeout (303\,s) is a cartesian explosion on a two-pattern MATCH without intermediate aggregation---a query optimization issue, not a data issue. Latencies for passing queries range from 1.3\,s to 19.4\,s.

\subsection{Public Health Trifecta (522K nodes, 546K edges)}

\begin{table}[H]
  \centering
  \caption{Public health benchmark results (40 queries).}
  \label{tab:ph-results}
  \begin{tabular}{@{}lrrrr@{}}
    \toprule
    \textbf{Category} & \textbf{Queries} & \textbf{Pass} & \textbf{Median (ms)} \\
    \midrule
    Health Determinants  & 20 & 20 & 27.4 \\
    Health Systems       & 10 & 10 & 12.2 \\
    Cross-KG / Cross-Indicator & 10 & 10 & 190.7 \\
    \midrule
    \textbf{Total}       & \textbf{40} & \textbf{40} & \textbf{15.1} \\
    \bottomrule
  \end{tabular}
\end{table}

Point lookups: 1--4\,ms. Single-hop traversals: 2--25\,ms. Aggregation scans: 10--134\,ms. Cross-KG joins: 66--478\,ms. Import time: 1.6\,s (Health Determinants, 286K nodes from \texttt{.sgsnap}) + 0.1\,s (Health Systems, 20K nodes).

\subsection{Cost Comparison}

\begin{table}[H]
  \centering
  \caption{Cost comparison for comparable billion-edge workload.}
  \label{tab:cost}
  \begin{tabular}{@{}lrl@{}}
    \toprule
    \textbf{System} & \textbf{Monthly Cost} & \textbf{Notes} \\
    \midrule
    Samyama (spot, load-and-query) & \$2.50 one-time & r6a.8xlarge, 3h46m \\
    Samyama (persistent)           & \$200--480/mo    & r8g.4xlarge--r6a.8xlarge \\
    Neo4j Aura Professional        & \$2,000+/mo      & Managed, 256GB tier \\
    AWS Neptune                    & \$1,500+/mo      & r5.8xlarge + storage \\
    TigerGraph Cloud               & \$5,000+/mo      & Enterprise tier \\
    \bottomrule
  \end{tabular}
\end{table}

% ============================================================
% 7. DISCUSSION & LIMITATIONS
% ============================================================
\section{Discussion \& Limitations}
\label{sec:discussion}

\paragraph{Failing queries.} Four biomedical queries fail: PM10 (article lacks grant data in snapshot), PM28 (cartesian explosion---timeout at 303\,s), CT18 (missing \texttt{CODED\_AS\_DRUG} edge type), DI05 (data gap between gene target and side effect sources). These are data completeness or query optimization issues, not engine limitations.

\paragraph{Snapshot portability.} Graphs loaded via \texttt{create\_edge\_stub()} bypass the Edge store, causing \texttt{export\_snapshot()} to omit edges. A v0.8.0 fix (edge store removal, adjacency-only serialization) will resolve this, enabling portable PubMed snapshots.

\paragraph{Single-node architecture.} All evaluations are single-node. A 3-node Raft cluster has been validated at smaller scale (426K nodes, all queries $<$1\,s), but distributed billion-edge federation is future work.

\paragraph{MVCC.} Read-snapshot isolation foundation is implemented; full MVCC is in progress. Concurrent read-write workloads on the billion-edge graph have not been evaluated.

\paragraph{Federation semantics.} Property-value bridge federation assumes that \texttt{iso\_code = 'IND'} in Health Determinants refers to the same entity as \texttt{iso\_code = 'IND'} in Health Systems. When KGs are imported as separate snapshots, Country nodes are duplicated (one per snapshot). Queries must join by property value, not node identity. A future entity resolution layer could merge duplicate nodes at import time.

% ============================================================
% 8. RELATED WORK
% ============================================================
\section{Related Work}
\label{sec:related}

\paragraph{Graph databases at scale.} Neo4j~supports billion-edge graphs in its Enterprise Edition with JVM-based storage and native clustering, but requires significant infrastructure investment. TigerGraph~achieves high throughput through pre-compiled GSQL and massively parallel processing but is proprietary. Kuzu~\citep{jin2023kuzu} and DuckPGQ~target embedded analytical workloads with columnar storage but lack server mode and distributed capabilities. Memgraph~provides in-memory C++ performance but without vector search or optimization solvers.

\paragraph{LDBC benchmarks.} The LDBC Graphalytics~\citep{iosup2016ldbc} and SNB Interactive benchmarks provide standardized evaluation of graph databases. Our work extends beyond LDBC synthetic datasets to real-world biomedical data at comparable scale.

\paragraph{Biomedical knowledge graphs.} Hetionet~\citep{himmelstein2017hetionet} integrates 47K nodes from 29 biomedical sources. PrimeKG~\citep{chandak2023primekg} provides 129K nodes for drug repurposing. Bio2RDF~\citep{belleau2008bio2rdf} uses RDF federation across SPARQL endpoints. Our work differs in scale (74.3M nodes vs.\ ${\sim}$100K), in using a single unified property graph (vs.\ RDF federation), and in providing a public health dimension absent from prior work.

\paragraph{Bulk loading.} The two-phase stub loading pattern draws on techniques from columnar databases~\citep{abadi2008column} (late materialization) and log-structured merge trees~\citep{oneil1996lsm} (append-only writes with background compaction). The CSR representation for adjacency lists is standard in graph analytics~\citep{iosup2016ldbc} but typically used read-only; our hybrid approach supports concurrent reads during compaction.

% ============================================================
% 9. CONCLUSION
% ============================================================
\section{Conclusion}

We have demonstrated that billion-edge cross-domain knowledge graph federation is achievable on commodity hardware with an open-source graph database. The key enablers are: (1)~two-phase bulk loading with stub nodes/edges achieving 13.6$\times$ memory reduction; (2)~hybrid CSR adjacency with mid-phase compaction enabling 1B edges on 256\,GB RAM; and (3)~property-value bridge federation requiring no ETL between knowledge graphs.

The 6-KG federation---spanning molecular biology (Pathways, Drug Interactions) through clinical medicine (Clinical Trials, PubMed) to population health (Disease Surveillance, Health Determinants, Health Systems)---achieves 136 of 140 benchmark queries with latencies from 1\,ms to 19\,s. The total compute cost is \$2.50.

\paragraph{Reproducibility.} All 8 knowledge graphs, Rust loaders, snapshots (\texttt{s3://samyama-data/snapshots/}), benchmark queries (CSV), and results are open-source. The Supabase \texttt{kg\_registry} table catalogs all KGs with node/edge counts, data sources, and S3 paths.

\paragraph{Future work.} Key directions include: portable snapshot export for stub-loaded graphs (v0.8.0); tiered hot/cold storage enabling billion edges on 32\,GB RAM (\$65/month); expanding the benchmark to 200 queries; MVCC for concurrent read-write workloads; and distributed federation across multiple Samyama nodes.

\paragraph{Availability.} Source code: \url{https://github.com/samyama-ai/samyama-graph}. Snapshots: \url{https://github.com/samyama-ai/samyama-graph/releases/tag/kg-snapshots-v6}. Documentation: \url{https://samyama-ai.github.io/samyama-graph-book/}.

% ============================================================
% REFERENCES
% ============================================================
\bibliographystyle{plainnat}
\bibliography{paper7_billion_edge}

\end{document}
